% Options for packages loaded elsewhere
% Options for packages loaded elsewhere
\PassOptionsToPackage{unicode}{hyperref}
\PassOptionsToPackage{hyphens}{url}
\PassOptionsToPackage{dvipsnames,svgnames,x11names}{xcolor}
%
\documentclass[
  11pt,
]{article}
\usepackage{xcolor}
\usepackage[margin=1in]{geometry}
\usepackage{amsmath,amssymb}
\setcounter{secnumdepth}{-\maxdimen} % remove section numbering
\usepackage{iftex}
\ifPDFTeX
  \usepackage[T1]{fontenc}
  \usepackage[utf8]{inputenc}
  \usepackage{textcomp} % provide euro and other symbols
\else % if luatex or xetex
  \usepackage{unicode-math} % this also loads fontspec
  \defaultfontfeatures{Scale=MatchLowercase}
  \defaultfontfeatures[\rmfamily]{Ligatures=TeX,Scale=1}
\fi
\usepackage{lmodern}
\ifPDFTeX\else
  % xetex/luatex font selection
\fi
% Use upquote if available, for straight quotes in verbatim environments
\IfFileExists{upquote.sty}{\usepackage{upquote}}{}
\IfFileExists{microtype.sty}{% use microtype if available
  \usepackage[]{microtype}
  \UseMicrotypeSet[protrusion]{basicmath} % disable protrusion for tt fonts
}{}
\makeatletter
\@ifundefined{KOMAClassName}{% if non-KOMA class
  \IfFileExists{parskip.sty}{%
    \usepackage{parskip}
  }{% else
    \setlength{\parindent}{0pt}
    \setlength{\parskip}{6pt plus 2pt minus 1pt}}
}{% if KOMA class
  \KOMAoptions{parskip=half}}
\makeatother
% Make \paragraph and \subparagraph free-standing
\makeatletter
\ifx\paragraph\undefined\else
  \let\oldparagraph\paragraph
  \renewcommand{\paragraph}{
    \@ifstar
      \xxxParagraphStar
      \xxxParagraphNoStar
  }
  \newcommand{\xxxParagraphStar}[1]{\oldparagraph*{#1}\mbox{}}
  \newcommand{\xxxParagraphNoStar}[1]{\oldparagraph{#1}\mbox{}}
\fi
\ifx\subparagraph\undefined\else
  \let\oldsubparagraph\subparagraph
  \renewcommand{\subparagraph}{
    \@ifstar
      \xxxSubParagraphStar
      \xxxSubParagraphNoStar
  }
  \newcommand{\xxxSubParagraphStar}[1]{\oldsubparagraph*{#1}\mbox{}}
  \newcommand{\xxxSubParagraphNoStar}[1]{\oldsubparagraph{#1}\mbox{}}
\fi
\makeatother


\usepackage{longtable,booktabs,array}
\usepackage{calc} % for calculating minipage widths
% Correct order of tables after \paragraph or \subparagraph
\usepackage{etoolbox}
\makeatletter
\patchcmd\longtable{\par}{\if@noskipsec\mbox{}\fi\par}{}{}
\makeatother
% Allow footnotes in longtable head/foot
\IfFileExists{footnotehyper.sty}{\usepackage{footnotehyper}}{\usepackage{footnote}}
\makesavenoteenv{longtable}
\usepackage{graphicx}
\makeatletter
\newsavebox\pandoc@box
\newcommand*\pandocbounded[1]{% scales image to fit in text height/width
  \sbox\pandoc@box{#1}%
  \Gscale@div\@tempa{\textheight}{\dimexpr\ht\pandoc@box+\dp\pandoc@box\relax}%
  \Gscale@div\@tempb{\linewidth}{\wd\pandoc@box}%
  \ifdim\@tempb\p@<\@tempa\p@\let\@tempa\@tempb\fi% select the smaller of both
  \ifdim\@tempa\p@<\p@\scalebox{\@tempa}{\usebox\pandoc@box}%
  \else\usebox{\pandoc@box}%
  \fi%
}
% Set default figure placement to htbp
\def\fps@figure{htbp}
\makeatother





\setlength{\emergencystretch}{3em} % prevent overfull lines

\providecommand{\tightlist}{%
  \setlength{\itemsep}{0pt}\setlength{\parskip}{0pt}}



 


\usepackage{amsmath,amssymb,enumitem,booktabs}
\setlist{itemsep=2pt,topsep=3pt}
\usepackage{fancyhdr}\pagestyle{fancy}\fancyhf{}
\fancyhead[L]{\small DS701 Fall 2026}\fancyhead[R]{\small Homework 4}
\fancyfoot[C]{\small\thepage}
\renewcommand{\headrulewidth}{0.4pt}
\usepackage[breakable]{tcolorbox}
% --- answer-space macros (problems PDF only; no-ops in the solutions PDF) ---
% \answerspace{2in}  vertical room for handwritten work, sized from the solution
% \blank             a short fill-in rule; \blank[1.5in] for a wider one
% \answerpage        start the next problem on a fresh page (Gradescope page matching)
\newcommand{\answerspace}[1]{\par\vspace*{#1}}  % \vspace* so room is not dropped at a page break
\newcommand{\blank}[1][1.2in]{\rule{#1}{0.4pt}}
\newcommand{\answerpage}{\clearpage}
\newtcolorbox{solutionbox}{breakable,colback=blue!4,colframe=blue!45!black,title=Solution,fonttitle=\bfseries,left=4pt,right=4pt,top=3pt,bottom=3pt}
\makeatletter
\@ifpackageloaded{caption}{}{\usepackage{caption}}
\AtBeginDocument{%
\ifdefined\contentsname
  \renewcommand*\contentsname{Table of contents}
\else
  \newcommand\contentsname{Table of contents}
\fi
\ifdefined\listfigurename
  \renewcommand*\listfigurename{List of Figures}
\else
  \newcommand\listfigurename{List of Figures}
\fi
\ifdefined\listtablename
  \renewcommand*\listtablename{List of Tables}
\else
  \newcommand\listtablename{List of Tables}
\fi
\ifdefined\figurename
  \renewcommand*\figurename{Figure}
\else
  \newcommand\figurename{Figure}
\fi
\ifdefined\tablename
  \renewcommand*\tablename{Table}
\else
  \newcommand\tablename{Table}
\fi
}
\@ifpackageloaded{float}{}{\usepackage{float}}
\floatstyle{ruled}
\@ifundefined{c@chapter}{\newfloat{codelisting}{h}{lop}}{\newfloat{codelisting}{h}{lop}[chapter]}
\floatname{codelisting}{Listing}
\newcommand*\listoflistings{\listof{codelisting}{List of Listings}}
\makeatother
\makeatletter
\makeatother
\makeatletter
\@ifpackageloaded{caption}{}{\usepackage{caption}}
\@ifpackageloaded{subcaption}{}{\usepackage{subcaption}}
\makeatother
\usepackage{bookmark}
\IfFileExists{xurl.sty}{\usepackage{xurl}}{} % add URL line breaks if available
\urlstyle{same}
\hypersetup{
  pdftitle={DS701 Homework 4: Probability, MLE, and EM},
  colorlinks=true,
  linkcolor={blue},
  filecolor={Maroon},
  citecolor={Blue},
  urlcolor={Blue},
  pdfcreator={LaTeX via pandoc}}


\title{DS701 Homework 4: Probability, MLE, and EM}
\usepackage{etoolbox}
\makeatletter
\providecommand{\subtitle}[1]{% add subtitle to \maketitle
  \apptocmd{\@title}{\par {\large #1 \par}}{}{}
}
\makeatother
\subtitle{Released Mon Sep 28 · due Mon Oct 5, 11:59 pm on Gradescope}
\author{}
\date{}
\begin{document}
\maketitle


\noindent\textbf{Name:}\ \rule{2.0in}{0.4pt}\qquad\textbf{BU ID:}\ \rule{1.2in}{0.4pt}
\vspace{6pt}

\textbf{How this works.} Answer \textbf{by hand on paper} (or a tablet
with a stylus), show your work, then scan and upload to Gradescope,
matching each problem to its pages. This is deliberate practice for the
closed-notes quizzes, which are in the same style --- so do it as much
as possible without a computer or an AI assistant. If you are unable to
do it without assistance, try AI assistance in ``socratic'' mode, where
you preface your questions by telling AI it is a tutor helping you learn
rather then give you answers right away. This homework is
\textbf{participation-graded}: a serious attempt at every problem earns
full credit, and worked solutions are released the day after the
deadline. Budget 1--2 hours.

\textbf{Covers:} Lecture 05 (Probabilistic Modeling) and Lecture 06
(Clustering III: GMM and EM). Bayes' rule for events, odds, and the
Bernoulli pmf are stated in the problems. Give numerical answers to 3
decimal places; keep the exact fractions when they are simple.

\subsection{Problem 1: A screening test, three
ways}\label{problem-1-a-screening-test-three-ways}

\emph{Practices: Bayes' rule with the base rate, and the odds /
likelihood-ratio form.}

A screening test for a disease \(D\) has \textbf{sensitivity}
\(P(+\mid D) = 0.90\) (positive in 90\% of people who have \(D\)) and
\textbf{specificity} \(P(-\mid \bar D) = 0.90\) (negative in 90\% of
people who do not). Bayes' rule for events, with the denominator
expanded by the law of total probability (with
\(C_k \in \{D, \bar D\}\)):
\[P(D\mid +) = \frac{P(+\mid D)\,P(D)}{P(+\mid D)\,P(D) + P(+\mid \bar D)\,P(\bar D)} .\]

\begin{enumerate}
\def\labelenumi{(\alph{enumi})}
\tightlist
\item
  In the general population the \textbf{prevalence} (base rate) is
  \(P(D) = 0.01\). A randomly chosen person tests positive. Compute
  \(P(D\mid +)\), showing the two terms of the denominator separately.
  In one sentence, why is the answer so far below the 90\% that ``the
  test is 90\% accurate'' suggests? \(P(D\mid +) =\) \blank[0.8in]
\end{enumerate}

\answerspace{2in}

\begin{enumerate}
\def\labelenumi{(\alph{enumi})}
\setcounter{enumi}{1}
\tightlist
\item
  The same test is used in a high-risk clinic where \(P(D) = 0.10\).
  Recompute \(P(D\mid +)\). Then state the practical lesson in one or
  two sentences: the test did not change --- what did, and what does
  that say about using one test both for mass screening and in a
  specialist clinic? \(P(D\mid +) =\) \blank[0.8in]
\end{enumerate}

\answerspace{2in}

\begin{enumerate}
\def\labelenumi{(\alph{enumi})}
\setcounter{enumi}{2}
\tightlist
\item
  \textbf{Odds form.} Define the odds of an event as
  \(O(D) = P(D)/P(\bar D)\) and the positive likelihood ratio as
  \(LR_+ = P(+\mid D)/P(+\mid \bar D)\). Show, starting from Bayes' rule
  written once for \(D\) and once for \(\bar D\), that
  \[O(D\mid +) = LR_+ \times O(D)\] (posterior odds = likelihood ratio ×
  prior odds --- no denominator needed). Optionally, then compute
  \(LR_+\) for this test, redo (a) and (b) in odds and convert back to
  probabilities to confirm your answers, and finally: a person from the
  general population tests positive on \textbf{two independent}
  administrations of the test --- what is \(P(D\mid ++)\)? (Independence
  here means the second test multiplies the odds by \(LR_+\) again.)
  \(LR_+ =\) \blank[0.6in]\quad \(P(D\mid ++) =\) \blank[0.8in]
\end{enumerate}

\answerspace{4in}

\answerpage

\subsection{Problem 2: Maximum likelihood, from
scratch}\label{problem-2-maximum-likelihood-from-scratch}

\emph{Practices: the MLE recipe --- likelihood, log, derivative, solve
--- on two models, and knowing why the log is there.}

\begin{enumerate}
\def\labelenumi{(\alph{enumi})}
\tightlist
\item
  \textbf{Bernoulli.} Each observation is \(x_n\in\{0,1\}\) with
  \(P(x_n = 1) = p\), so \(P(x_n\mid p) = p^{x_n}(1-p)^{1-x_n}\). For
  \(N\) i.i.d. observations with \(k = \sum_n x_n\) ones: write the
  likelihood \(L(p)\), the log-likelihood \(\ell(p)\), set
  \(d\ell/dp = 0\), and solve for \(\hat p\). Then evaluate \(\hat p\)
  for the data \(1,0,1,1,0,1,1,1,0,1\). (Say in one line why the
  critical point is a maximum and not a minimum.) For the data,
  \(\hat p =\) \blank[0.8in]
\end{enumerate}

\answerspace{3in}

\begin{enumerate}
\def\labelenumi{(\alph{enumi})}
\setcounter{enumi}{1}
\tightlist
\item
  \textbf{Gaussian mean and variance.} Start from the Gaussian
  log-likelihood
  \[\ell(\mu,\sigma^2) = -\frac N2\log(2\pi\sigma^2) - \frac{1}{2\sigma^2}\sum_{n=1}^N (x_n-\mu)^2 .\]
  Derive \(\hat\mu\) from \(\partial\ell/\partial\mu = 0\). Then treat
  \(v = \sigma^2\) as the variable, derive \(\hat\sigma^2\) from
  \(\partial\ell/\partial v = 0\) (substituting \(\hat\mu\)), and
  evaluate both on the data \(2, 3, 5, 6, 9\). Compare \(\hat\sigma^2\)
  with the sample variance \(s^2\) that divides by \(N-1\) --- which is
  larger, and by what factor here? \(\hat\mu =\)
  \blank[0.6in]\quad \(\hat\sigma^2 =\) \blank[0.6in]\quad \(s^2 =\)
  \blank[0.6in]\quad factor: \blank[0.6in]
\end{enumerate}

\answerspace{3in}

\begin{enumerate}
\def\labelenumi{(\alph{enumi})}
\setcounter{enumi}{2}
\tightlist
\item
  \textbf{Why the log?} Give the two reasons the lecture gives for
  maximizing \(\ell\) rather than \(L\): one about \emph{where} the
  maximum is, one about the \emph{form} of the function. Then add the
  practical one: for \(N = 2000\) observations from (a), every factor of
  \(L\) is at most \(1\) --- say roughly how small \(L\) is at best
  (e.g.~compare \(0.5^{2000}\) with the smallest positive double, about
  \(10^{-308}\)) and why \(\ell\) does not have this problem.
\end{enumerate}

\answerspace{2.2in}

\answerpage

\subsection{Problem 3: One EM iteration by
hand}\label{problem-3-one-em-iteration-by-hand}

\emph{Practices: the E-step and M-step formulas on numbers small enough
to see what they do, and the argument for why EM behaves.}

A one-dimensional mixture of \(K=2\) Gaussians, both with variance
\(\sigma^2 = 1\) (\textbf{keep the variances fixed at 1 throughout} ---
update only means and weights). Current parameters: \(\mu_1 = 0\),
\(\mu_2 = 2\), \(\pi_1 = \pi_2 = 0.5\). Data: \(x_1 = 0\), \(x_2 = 1\),
\(x_3 = 2\), \(x_4 = 4\). Useful values: \(e^{-0.5}\approx0.607\),
\(e^{-2}\approx0.135\), \(e^{-8}\approx0.00034\); note that with equal
variances the factor \(1/\sqrt{2\pi}\) is the same in every term and
cancels in the responsibilities.

\begin{enumerate}
\def\labelenumi{(\alph{enumi})}
\tightlist
\item
  \textbf{E-step.} For each data point compute the responsibilities
  \(\gamma_{i1}, \gamma_{i2}\) (3 decimals). Before computing, say for
  which point the answer must be exactly \(0.5\) and why.
  \((\gamma_{i1},\gamma_{i2})\) for \(x_1,\dots,x_4\):
  \blank[1in] \blank[1in] \blank[1in] \blank[1in]
\end{enumerate}

\answerspace{2.5in}

\begin{enumerate}
\def\labelenumi{(\alph{enumi})}
\setcounter{enumi}{1}
\tightlist
\item
  \textbf{M-step.} Compute \(N_1, N_2\), the new means
  \(\mu_1', \mu_2'\) and the new weights \(\pi_1', \pi_2'\). In one
  sentence each: why did \(\mu_2\) move more than \(\mu_1\), and why did
  the weights become unequal? \(N_1 =\) \blank[0.5in] \(N_2 =\)
  \blank[0.5in] \(\mu_1' =\) \blank[0.5in] \(\mu_2' =\)
  \blank[0.5in] \(\pi_1' =\) \blank[0.5in] \(\pi_2' =\) \blank[0.5in]
\end{enumerate}

\answerspace{3.4in}

\begin{enumerate}
\def\labelenumi{(\alph{enumi})}
\setcounter{enumi}{2}
\tightlist
\item
  The lecture states that the log-likelihood \(\ell(\theta)\) never
  decreases from one EM iteration to the next. Explain the mechanism in
  words: what the E-step computes, what the M-step maximizes and why
  that has a closed form, and why alternating the two cannot lower
  \(\ell\). Then say what this guarantees about the sequence
  \(\ell^{(t)}\), what it does \textbf{not} guarantee, and what is done
  about that in practice (relate to Homework 3, Problem 1).
\end{enumerate}

\answerspace{2.8in}

\begin{enumerate}
\def\labelenumi{(\alph{enumi})}
\setcounter{enumi}{3}
\tightlist
\item
  \textbf{Hard vs.~soft.} Without recomputing, describe what the four
  responsibilities in (a) would look like if both variances were
  \(\sigma^2 = 0.01\) instead of \(1\), and what the M-step for the
  means would then reduce to. Which algorithm is that? State the one
  practical difference and the one conceptual difference between GMM/EM
  and \(k\)-means that the lecture draws.
\end{enumerate}

\answerspace{2.2in}




\end{document}
